\clearpage\begingroup\NoSourceHead\SourceLeaf{274}
\providecommand{\ObsSec}[2]{#1\mathord{\overset{\mathrm{s}}{,}}#2}
\providecommand{\ObsArc}[2]{#1\mathord{\overset{\prime\prime}{,}}#2}
\providecommand{\ObsJupiter}{\mathord{\tikz[baseline=-.1ex,x=1em,y=1em,line width=.045em]{
  \draw (.02,.48) .. controls (.04,.70) and (.33,.69) .. (.31,.46) .. controls (.29,.29) and (.14,.19) .. (.02,.13)--(.49,.13);
  \draw (.37,.64)--(.37,-.06);}}}
\providecommand{\ObsVesta}{\mathord{\tikz[baseline=-.1ex,x=1em,y=1em,line width=.045em]{
  \draw (.04,.02) rectangle (.40,.38);
  \draw (.22,.38)--(.22,.48) (.12,.63)--(.22,.48)--(.32,.63);}}}

\begin{center}
{\large [VERSCHIEDENE BEOBACHTUNGEN\\IN GÖTTINGEN.]}\\[1em]
NACHLASS, VERÖFFENTLICHUNGEN UND\\BRIEFWECHSEL.
\end{center}
\begin{center}\small
\rule{.92\linewidth}{.35pt}\\
Tagebuch der astronomischen Beobachtungen auf der königl. Sternwarte zu Göttingen.\\
{}[Erster Band; Angefangen] November 1808.\\
\rule{.92\linewidth}{.35pt}
\end{center}
\noindent [1. Seite] 12, 1808 December 18.

\{Austritt des I. $\ObsJupiter$-Trabanten, beobachtet von Herrn Prof[essor] \textsc{Gauss} mit 10 f. Reflector, schwächste Vergr.
\begin{center}\small\setlength{\tabcolsep}{3pt}
\begin{tabular}{l lrrr}
\raisebox{-.5\normalbaselineskip}{um $22^{\mathrm h}29^{\mathrm m}58^{\mathrm s}$ nach Shelton\}\textsuperscript{[*)]} ist in}&Sternzeit&$22^{\mathrm h}$&$31^{\mathrm m}$&$\ObsSec80$\\
&Mittl. Zeit&4&42&$44{,}7$\textsuperscript{[**)]}.
\end{tabular}
\end{center}
\begingroup\renewcommand{\thefootnote}{[*)}\footnotetext{\textsc{Hardings} Handschrift.]}\endgroup
\begingroup\renewcommand{\thefootnote}{[**)}\footnotetext{\textsc{Gauss}’ Handschrift.]}\endgroup
\noindent [2. Seite] 55, 1809 September 4.

Der Eintritt von $\lambda$ Geminorum geschah unter dem Horizont; der Austritt wurde beobachtet am dunkeln aber sichtbaren Mondsrande
\begin{center}\small\setlength{\tabcolsep}{3pt}\renewcommand{\arraystretch}{1.2}
\begin{tabular}{l r|l}
von Herrn Dr. \textsc{Schumacher} mit&&also in $\ast$-Zeit $(1^{\mathrm m}\ObsSec{47}{18}+\text{Uhrzeit})$\\
\quad Short&$0^{\mathrm h}28^{\mathrm m}\ObsSec{13}3$&$0^{\mathrm h}30^{\mathrm m}\ObsSec0{48}=13^{\mathrm h}35^{\mathrm m}\ObsSec86$ M. Z.\\
von mir mit Herschel&$28\;13{,}5$&$30\;0{,}68=13\;35\;8{,}8$\\
von Herrn Professor \textsc{Harding}&&\\
\quad mit gr. Dollond&$28\;14{,}0$&$30\;1{,}18=13\;35\;9{,}3$.
\end{tabular}
\end{center}
\noindent\hspace*{\parindent}Sehr gute Beobachtung\textsuperscript{[***)]}.
\begingroup\renewcommand{\thefootnote}{[***)}\footnotetext{Vgl. auch Brief von \textsc{Gauss} an \textsc{Olbers} vom 12. September 1809 und an \textsc{Bessel} vom gleichen Tage.]}\endgroup
\clearpage\SourceLeaf{275}\HistoricalHead{270}{L}{PRAKTISCHE UND SPHÄRISCHE ASTRONOMIE. NACHLASS.}
\noindent [3. Seite] 58, 1809 October 3.

Bei zieml[ich] heitrer Luft, wiewohl späterhin ein starker Nebel sich verbreitete, wurde die Vesta viermal mit einem Sterne der \textit{Histoire Céleste} $a$, p. 253, 1797 Febr. 17
\begin{center}\begin{tabular}{c|c|c|c|c}
8&$6^{\mathrm h}56^{\mathrm m}\ObsSec{33}5$&$6^{\mathrm h}57^{\mathrm m}\ObsSec24$&—&$28^0\,57'\,20''$
\end{tabular}\end{center}
\noindent verglichen; ein anderer anonymer Stern, $b$, [der] etwa $15'$ nördlicher stehen mochte, wurde mitbeobachtet. Vesta hatte die 8. Grösse
\begin{center}\small\setlength{\tabcolsep}{3pt}\renewcommand{\arraystretch}{1.2}
\begin{tabular}{r r l|r r|r r}
\multicolumn{3}{c}{$\ObsVesta$}&\multicolumn{2}{c}{$a$}&\multicolumn{2}{c}{$b$}\\\hline
$2^{\mathrm h}\;1^{\mathrm m}53^{\mathrm s}$&$4^{\mathrm m}37^{\mathrm s}$&S[üdl.]&$4^{\mathrm m}47^{\mathrm s}\mathord{:}\mathord{:}\mathord{:}$&$7^{\mathrm m}25^{\mathrm s}$&$6^{\mathrm m}\;\ObsSec55$&$8^{\mathrm m}\ObsSec{53}5$\\
$13\quad2$&$16\quad1$&&$15\quad50$&$18\quad50{,}5$&$17\quad51$&$19\quad38$\\
$29\quad19$&$30\quad36$&S&$32\quad16$&$33\quad17$&$32\quad41$&$35\quad38$\\
$40\quad25{,}5$&$43\quad23$&&$43\quad14$&$46\quad11$&$44\quad52$&$47\quad21$\\
$55\quad11$&$56\quad53{,}5$&N&$57\quad52$&$59\quad46$&—&—
\end{tabular}
\end{center}
\noindent Reduction des Sterns $a$ aus der \textit{Histoire Céleste} . . . $[\alpha]$ appar. $104^0\,29'\,\ObsArc3{82}$ $[\delta\ \text{app.}]$ $19^0\,50'\,\ObsArc{17}{22}$.
\begin{center}\small
\begin{minipage}[t]{.49\linewidth}\centering
Es ging also $\ObsVesta$ dem $a$ vor:\\[.4em]
$\left.\begin{array}{lrrr}
\text{Uhrzeit}&2^{\mathrm h}&28^{\mathrm m}&\ObsSec{12}9\\
\ast\ \mathrm Z.&2&33&16{,}9\\
\mathrm{M.\ Z.}&13&44&3
\end{array}\right\}\ \text{um}\ 2^{\mathrm m}\ObsSec{48}5$\\[.5em]
$[\alpha=]\ 103^0\,46'\,\ObsArc{56}3$
\end{minipage}\hfill
\begin{minipage}[t]{.49\linewidth}\centering
und war nördlicher als $a$:\\[.4em]
Uhrzeit $2^{\mathrm h}42^{\mathrm m}\ObsSec{59}9:$ [um] $\ObsArc{54}1$\\[3.3em]
$[\delta=]\ 19^0\,51'\,\ObsArc{11}3$.
\end{minipage}
\end{center}
\noindent\hspace*{\parindent}Nach der ersten Beob. der Stern $b$ $16'12''$ N[ördlich] von der $\ObsVesta$; nach 2. und 4. [Beobachtung] folgt $b$ $1^{\mathrm m}\ObsSec{24}1$ dem $a$.
\medskip\noindent [4. Seite] 60, 1809 October 4.

Vergleichung der Vesta mit dem Stern $a$ von gestern
\begin{center}\small\setlength{\tabcolsep}{3pt}\renewcommand{\arraystretch}{1.2}
\begin{tabular}{r r l|r r l}
\multicolumn{3}{c}{$\ObsVesta$}&\multicolumn{2}{c}{$a$}&\\\cline{1-5}
$1^{\mathrm h}55^{\mathrm m}23^{\mathrm s}$&$58^{\mathrm m}20^{\mathrm s}$&N[ördlich]&$57^{\mathrm m}13^{\mathrm s}$&$60^{\mathrm m}10^{\mathrm s}$&\\
$2\quad2\quad8$&$2\quad49$&&$3\quad55$&$4\quad43$&\\
$6\quad30$&$9\quad29$&&$8\quad20$&$11\quad18{,}5$&\\
$15\quad48$&$16\quad43$&S&$17\quad40$&$18\quad30$&\\
$30\quad16$&$32\quad19$&N&$32\quad4$&$34\quad11$&[Beob.] Dr. \textsc{Schumacher}.
\end{tabular}
\end{center}
\clearpage\SourceLeaf{276}\HistoricalHead{271}{R}{VERSCHIEDENE BEOBACHTUNGEN IN GÖTTINGEN.}
\begin{center}\small
\begin{minipage}[t]{.49\linewidth}\centering
Also $\ObsVesta$ geht $a$ vor\\[.4em]
\setlength{\tabcolsep}{2pt}\begin{tabular}{lrrrl}
Uhrzeit&$2^{\mathrm h}$&$2^{\mathrm m}$&$\ObsSec{25}5$&[um] $1^{\mathrm m}\ObsSec{49}9$\\
$\ast$ Zeit&2&7&$32{,}0$&\\
M. Z.&13&14&27&
\end{tabular}\\[.5em]
$[\alpha=]\ 104^0\,1'\,\ObsArc{35}6$
\end{minipage}\hfill
\begin{minipage}[t]{.49\linewidth}\centering
$\ObsVesta$ ist nördlicher als $a$:\\[.4em]
\setlength{\tabcolsep}{2pt}\begin{tabular}{lrrrl}
{}[Uhrzeit]&$2^{\mathrm h}$&$9^{\mathrm m}$&$\ObsSec{22}0$&[um] $\ObsArc{11}7$\\
{}[$\ast$ Zeit]&2&14&$28{,}5$&
\end{tabular}\\[1.7em]
$[\delta=]\ 19^0\,50'\,\ObsArc{28}9$
\end{minipage}
\end{center}
\noindent\hspace*{\parindent}Mira Ceti heute Abend 5. 6. Grösse.
\clearpage\endgroup
